\chapter{Getting Data In and Out}
\label{ch:data}

\section{Supported formats}
\label{sec:formats}

All structure I/O goes through \texttt{ase.io}, so Calango reads and writes
every format ASE knows. The file dialog offers the common ones directly:

\begin{center}
\small
\begin{tabular}{@{}lp{9.2cm}@{}}
\toprule
Family & Extensions \\
\midrule
XYZ            & \file{.xyz}, \file{.extxyz} \\
Crystallography& \file{.cif} \\
VASP           & \file{POSCAR}, \file{CONTCAR}, \file{.vasp} \\
ASE            & \file{.traj} \\
Quantum ESPRESSO & \file{.in}, \file{.pwi}, \file{.pwo}, \file{.out} \\
CASTEP         & \file{.cell} \\
LAMMPS         & \file{.data}, \file{.dump}, \file{.lammpstrj} \\
Gaussian       & \file{.gjf}, \file{.com} \\
SHELX          & \file{.res} \\
\bottomrule
\end{tabular}
\end{center}

\begin{calnote}
Extended-XYZ per-atom arrays are imported as \emph{scalar fields} and
\emph{vector fields}. A \texttt{charges} or \texttt{forces} column becomes
selectable in the \ui{Custom property} colour mode, and \texttt{forces} or
\texttt{momenta} enable the force and velocity arrow overlays. This is the
simplest route for visualising per-atom data computed elsewhere: write it
as an extxyz column and open the file.
\end{calnote}

\section{Opening structures and trajectories}

\begin{description}
  \item[\menu{File}\menusep\menu{Open}\menusep\menu{Structure}]
    \kbd{Ctrl}+\kbd{O}.
  \item[\menu{File}\menusep\menu{Open}\menusep\menu{Trajectory}]
    \kbd{Ctrl}+\kbd{T}.
\end{description}

Both entries read \emph{every} frame in the file. A single-frame file opens
as a plain structure tab; a multi-frame file opens as a trajectory with the
timeline active. In practice the two entries differ only in the file
filter, so either works for either kind of file.

Opening a \file{.calproj} file --- from the command line, your file
manager, or the structure dialog --- restores the saved session rather than
loading a geometry.

\section{Saving structures and trajectories}

\begin{description}
  \item[\menu{File}\menusep\menu{Save}\menusep\menu{Structure As}]
    \kbd{Ctrl}+\kbd{Shift}+\kbd{S}. Writes the currently displayed frame.
    The name is suggested from the document's own, with the default
    extension; the format follows the filter you pick, and a name typed
    without an extension is given the filter's own.
  \item[\menu{File}\menusep\menu{Save}\menusep\menu{Trajectory As}]
    \kbd{Ctrl}+\kbd{Shift}+\kbd{T}. Writes all frames of the active
    document as extended XYZ, multi-frame XYZ, ASE \file{.traj}, or
    multi-model PDB.
\end{description}

\begin{calnote}
\textbf{Extended XYZ is the default everywhere.} It is pre-selected in every
open and save dialog in the application --- structures, trajectories, NEB
endpoints, dataset imports --- and it is what \menu{Save Structure As}
writes unless you choose otherwise. The reason is that it is the only format
in the list that round-trips everything a Calango document carries: cell and
periodicity, per-atom magnetic moments and charges, forces and velocities,
and arbitrary extra columns. Plain \file{.xyz} silently drops the cell; CIF
drops the calculator results; POSCAR drops nearly everything but the
positions.

Every other format is still one entry down in the same list, so nothing has
become harder to reach --- only the \emph{default} changed, to the format
that loses nothing.
\end{calnote}

\section{The database and preset browser}
\label{sec:database}

\menu{Build}\menusep\menu{From Database} opens a two-tab browser.

\subsection{Presets}

Seven benchmark structures ship with Calango, each annotated with a
suggested calculator:

\begin{center}
\small
\begin{tabular}{@{}lll@{}}
\toprule
Preset & File & Suggested calculator \\
\midrule
Diamond (bulk C)       & \file{diamond.vasp}            & MACE-MP-0 \\
MoS$_2$ 2H (bulk)      & \file{mos2\_2h\_bulk.vasp}     & MACE-MP-0 \\
Graphene monolayer     & \file{graphene.vasp}           & MACE-MP-0 \\
MoS$_2$ 1H (monolayer) & \file{mos2\_1h\_monolayer.vasp}& MACE-MP-0 \\
Benzene                & \file{benzene.xyz}             & MACE-OFF (or EMT) \\
Naphthalene            & \file{naphthalene.xyz}         & MACE-OFF \\
Coronene               & \file{coronene.xyz}            & MACE-OFF \\
\bottomrule
\end{tabular}
\end{center}

Select a row and press \ui{Load Selected}, or double-click it.

\subsection{Materials Project}

The second tab fetches structures by Materials Project ID.

\begin{description}
  \item[\ui{API key}] Masked field, pre-filled from your saved setting or
    from \opt{MP\_API\_KEY} in the environment. It is stored as you type.
  \item[\ui{Material ID}] For example \texttt{mp-149} (silicon). Pressing
    \kbd{Enter} triggers the fetch.
  \item[\ui{.env file}] Shows the current environment file path, with
    \ui{Browse\ldots} to point elsewhere and \ui{Reload} to re-import the
    key (Section~\ref{sec:envfile}).
  \item[\ui{Fetch Structure}] Downloads the entry and opens it in a new
    tab, reporting the formula and atom count.
\end{description}

\begin{calwarn}
Needs a network connection and a personal API key from
\url{https://materialsproject.org/api}.
\end{calwarn}
